\documentclass[preprint,12pt]{elsarticle}
\journal{Ecological Modelling}
\date{4 November 2026}
\begin{document}
\begin{frontmatter}
\title{Energy flows in a small estuary}
\author[a]{Ada Lovelace\corref{cor1}}
\ead{ada@example.org}
\author[a,b]{Howard T. Odum}
\cortext[cor1]{Corresponding author}
\address[a]{Department of Systems Ecology, University of Florida, Gainesville, USA}
\address[b]{Marine Science Institute, Port Aransas, USA}
\begin{abstract}
We measure the energy flows of a small estuary over one year and compare
them with a simple model. The model reproduces the seasonal pattern.
\end{abstract}
\begin{keyword}
emergy \sep estuary \sep systems ecology
\end{keyword}
\end{frontmatter}

\section{Introduction}
Estuaries collect energy from the tide, the river and the sun. This paper
follows the method of Odum \cite{odum_systems_1983}.

\subsection{Study site}
The estuary is small enough to be sampled in a day, which lets every flow
be measured on the same day of each month of the year.

\section{Results}
The model reproduces the seasonal pattern of the measured flows.

\bibliographystyle{elsarticle-num}
\bibliography{refs}
\end{document}
